% SHARD-ID: AQM-36-NILPOTENT-THICKENING-WEYL-FIELDS
% SHARD-TITLE: Nilpotent Thickening Weyl Fields
% SHARD-SUMMARY: Defines the nilpotent-sensitive Weyl layer for \(\mathbb F_3[t]/(\pi^e)\).
% SHARD-SUMMARY: Proves the top-coefficient character is generating and gives an irreducible Weyl representation.
% SHARD-SUMMARY: Interprets thickenings as jet degrees of freedom rather than additional closed points.
% SHARD-KEYWORDS: nilpotent thickening, finite chain ring, Frobenius ring, Weyl-Heisenberg, jet, Artin ring

\section{Nilpotent Thickening Weyl Fields}
\label{sec:nilpotent-thickening-weyl-fields}

\subsection{Status and Source Anchors}

This shard adds the nilpotent-sensitive layer proposed after AQM-34 and
AQM-35.  The finite Frobenius-ring source anchors are
\path{references/heisenberg_weil/gluesing_luerssen_pllaha_1710_09884_source/StabCodesFrob5.tex},
lines 224--249 for Frobenius rings and generating characters, lines
286--323 for \(X(a)\) and \(Z(b)\), lines 353--355 for the multiplication
law, and lines 436--452 for the error-basis statement.  Sidana--Kashyap,
\path{references/heisenberg_weil/sidana_kashyap_2202_00248_source/EAQECCs_over_rings_v15.tex},
lines 103--110 and 180--197, gives the same local Frobenius-ring Pauli
framework.  The trace pairing input for residue fields is the Stacks Project
snapshot \path{references/algebraic_geometry/stacks_project_algebra.tex},
lines 45278--45284.  The explicit generating-character proof below is local
to this shard.

\subsection{Lamport Dependency Skeleton}

\[
\begin{array}{rcl}
D1&:&k=\mathbb F_3,\ \pi\in k[t]\text{ monic irreducible},\
      e\ge1,\ A=k[t]/(\pi^e).\\
D2&:&\mathfrak m=(\pi)/(\pi^e),\ K=A/\mathfrak m=k[t]/(\pi),\
      d=\deg\pi.\\
L1&:&A\text{ has unique }\pi\text{-adic expansions and ideals }
      \mathfrak m^r.\\
D3&:&\lambda(a)=\Tr_{K/k}(a_{e-1})\text{ from the top }\pi\text{-coefficient},
      \quad \chi(a)=\exp(2\pi i\lambda(a)/3).\\
L2&:&\chi\text{ is generating: }b\ne0\Rightarrow
      (a\mapsto\chi(ba))\text{ is nontrivial}.\\
D4&:&\mathcal H_A=\ell^2(A),\quad E_A=A^2,\quad
      W_A(q,p)=T(q)R(p).\\
T1&:&W_A\text{ is a projective unitary Weyl representation with multiplier }
      \chi(pq').\\
T2&:&\{W_A(q,p):(q,p)\in A^2\}\text{ is an operator basis of }
      \operatorname{End}_\mathbb C(\mathcal H_A).\\
T3&:&\mathfrak m^r\text{ position jets pair with }
      \mathfrak m^{e-1-r}\text{ momentum jets in the associated graded}.\\
T4&:&\text{For }R_p=\mathbb F_3[t]/(p),\text{ the thickened layer has }
      \dim\mathcal H=3^{\deg p}.
\end{array}
\]

\subsection{The Local Artin Ring L1--L2}

Let
\[
  A=A_{\pi,e}=k[t]/(\pi^e),\qquad
  \mathfrak m=(\pi)/(\pi^e),\qquad K=k[t]/(\pi).
\]
Every class \(a\in A\) has a unique expansion
\[
  a=a_0+a_1\pi+\cdots+a_{e-1}\pi^{e-1},
\]
where each \(a_j\) is represented by a polynomial of degree \(<d=\deg\pi\).
This follows by repeated Euclidean division by the monic polynomial \(\pi\).

\paragraph{Lemma \(L1\).}
The ideals of \(A\) are exactly
\[
  A=\mathfrak m^0\supset\mathfrak m\supset\cdots
  \supset\mathfrak m^{e-1}\supset\mathfrak m^e=0.
\]
Every nonzero \(a\in A\) has a unique valuation \(v(a)=r\), \(0\le r<e\),
with \(a=\pi^r u\) for a unit \(u\in A^\times\).

\emph{Proof.}
If \(a\ne0\), take the least \(r\) for which the coefficient \(a_r\) is
nonzero.  Then \(a=\pi^r u\), where the residue of \(u\) in \(K\) is
\(a_r\ne0\), so \(u\) is a unit.  Let \(I\subset A\) be a nonzero ideal and
choose an element of minimal valuation \(r\).  Multiplying by the inverse of
its unit part gives \(\pi^r\in I\), hence \(\mathfrak m^r\subset I\).
Minimality of \(r\) gives \(I\subset\mathfrak m^r\), so \(I=\mathfrak m^r\).

Define
\[
  \lambda(a)=\Tr_{K/k}(a_{e-1}),\qquad
  \chi(a)=\exp\!\left(\frac{2\pi i}{3}\lambda(a)\right).
\]
Here \(a_{e-1}\in K\) is the top \(\pi\)-coefficient.  For \(e=1\) this is
the residue-field character used in AQM-34.

\paragraph{Lemma \(L2\).}
If \(b\ne0\) in \(A\), then the additive character
\[
  a\longmapsto\chi(ba)
\]
is nontrivial.  Hence the map \(b\mapsto(a\mapsto\chi(ba))\) identifies \(A\)
with the additive character group \(\widehat A\).

\emph{Proof.}
Write \(b=\pi^r u\) with \(u\in A^\times\).  Let \(\bar u\) be the residue of
\(u\) in \(K\).  Since \(K\) is finite, Frobenius \(x\mapsto x^3\) is
injective and hence surjective; thus \(K\) is perfect and \(K/k\) is finite
separable.  The cited Stacks trace-pairing source says
\((x,y)\mapsto\Tr_{K/k}(xy)\) is nondegenerate.
Choose \(\bar c\in K\) with \(\Tr_{K/k}(\bar u\bar c)\ne0\), and lift it to
\(c\in A\) of \(\pi\)-degree \(0\).  Set
\[
  a=\pi^{e-1-r}c.
\]
Then \(ba=\pi^{e-1}uc\) modulo \(\pi^e\), so the top coefficient of \(ba\) is
\(\bar u\bar c\).  Therefore \(\chi(ba)\ne1\).  The map \(A\to\widehat A\) is
injective, and both finite groups have cardinality \(|A|\), so it is an
isomorphism.

\subsection{The Thickened Weyl Representation T1--T2}

Let
\[
  \mathcal H_A=\ell^2(A)
\]
with basis \(|s\rangle\), \(s\in A\).  For \(q,p\in A\), define
\[
  T(q)|s\rangle=|s+q\rangle,\qquad
  R(p)|s\rangle=\chi(ps)|s\rangle,\qquad
  W_A(q,p)=T(q)R(p).
\]

\paragraph{Theorem \(T1\).}
For \(q,p,q',p'\in A\),
\[
  W_A(q,p)W_A(q',p')
  =
  \chi(pq')W_A(q+q',p+p')
\]
and
\[
  W_A(q,p)W_A(q',p')
  =
  \chi(pq'-p'q)W_A(q',p')W_A(q,p).
\]
Thus \(\mu_3\times A^2\), with multiplication
\[
  (\zeta,q,p)(\eta,q',p')
  =
  (\zeta\eta\chi(pq'),q+q',p+p'),
\]
acts unitarily on \(\mathcal H_A\) by
\((\zeta,q,p)\mapsto\zeta W_A(q,p)\).

\emph{Proof.}
\(T(q)\) is a permutation unitary and \(R(p)\) is diagonal unitary.  Also
\[
  R(p)T(q')|s\rangle
  =
  \chi(p(s+q'))|s+q'\rangle
  =
  \chi(pq')T(q')R(p)|s\rangle .
\]
Substituting this identity into \(T(q)R(p)T(q')R(p')\) gives the product law.
The commutator law follows by comparing with the same product law in reversed
order.

\paragraph{Theorem \(T2\).}
\[
  \operatorname{span}_{\mathbb C}\{W_A(q,p):(q,p)\in A^2\}
  =
  \operatorname{End}_{\mathbb C}(\mathcal H_A).
\]
Consequently the thickened Weyl representation is irreducible.

\emph{Proof.}
Use the Hilbert--Schmidt inner product on \(\operatorname{End}(\mathcal H_A)\).  If
\(q\ne q'\), then \(T(q'-q)\) has no fixed basis vector, so
\(\Tr(W_A(q,p)^*W_A(q',p'))=0\).  If \(q=q'\) but \(p\ne p'\), then the trace
is
\[
  \sum_{s\in A}\chi((p'-p)s),
\]
which is \(0\) because \(s\mapsto\chi((p'-p)s)\) is a nontrivial character by
L2.  Thus the \(|A|^2\) Weyl operators are pairwise orthogonal and nonzero.
Since \(\dim_\mathbb C\operatorname{End}(\ell^2(A))=|A|^2\), they form a basis.  Any
invariant subspace is invariant under this full matrix algebra, hence is
\(0\) or \(\mathcal H_A\).

\subsection{Jet Interpretation T3}

The filtration
\[
  A\supset\mathfrak m\supset\cdots\supset\mathfrak m^{e-1}\supset0
\]
is the algebraic jet filtration.  Its associated graded ring is
\[
  \operatorname{gr}_{\mathfrak m}A
  \cong K[\varepsilon]/(\varepsilon^e),
\]
with \(\mathfrak m^r/\mathfrak m^{r+1}\cong K\varepsilon^r\).

\paragraph{Lemma \(L3\).}
For \(0\le r\le e\), the annihilator of \(\mathfrak m^r\) under the pairing
\[
  A\times A\longrightarrow\mu_3,\qquad (x,y)\longmapsto\chi(xy)
\]
is \(\mathfrak m^{e-r}\).

\emph{Proof.}
If \(y\in\mathfrak m^{e-r}\), then \(xy\in\mathfrak m^e=0\) for all
\(x\in\mathfrak m^r\), so \(y\) annihilates \(\mathfrak m^r\).  Conversely,
if \(y\notin\mathfrak m^{e-r}\), write \(y=\pi^s u\) with \(s<e-r\).  Then
\(x=\pi^{e-1-s}c\) lies in \(\mathfrak m^r\) for any \(c\in A\) of degree
\(0\).  The proof of L2 chooses \(c\) so that \(\chi(xy)\ne1\).  Hence \(y\)
does not annihilate \(\mathfrak m^r\).

Thus a position jet in \(\mathfrak m^r/\mathfrak m^{r+1}\) pairs, in the
associated graded, with a momentum jet in
\(\mathfrak m^{e-1-r}/\mathfrak m^{e-r}\).  In particular, the ordinary
residue value \(A/\mathfrak m\) is dual to the socle momentum layer
\(\mathfrak m^{e-1}\), not to a new closed point.

The old residue qutrit embeds as the zeroth-position/socle-momentum subsystem:
choose lifts \(\tilde q,\tilde p\in A\) of \(q,p\in K\), and use
\[
  q\longmapsto \tilde q,\qquad p\longmapsto \pi^{e-1}\tilde p.
\]
On the subspace spanned by basis vectors with only degree-\(0\) coefficient,
these operators reproduce the residue-field Weyl law.

\subsection{Global Polynomial Quotients T4}

For
\[
  p=\prod_{i=1}^r\pi_i^{e_i},
\]
set \(A_i=\mathbb F_3[t]/(\pi_i^{e_i})\).  By Chinese remainders,
\[
  R_p=\mathbb F_3[t]/(p)\cong\prod_iA_i.
\]
The thickened arithmetic field is
\[
  \mathcal H_p^{\rm thick}
  =
  \bigotimes_i\ell^2(A_i)
  \cong \ell^2(R_p),
\qquad
  E_p^{\rm thick}=\bigoplus_i A_i^2\cong R_p^2.
\]
Its Weyl operators are tensor products of the local \(W_{A_i}\), equivalently
translations and phase multiplications on \(\ell^2(R_p)\) using the product
character.  The operator-basis theorem tensors over \(i\), so
\[
  \operatorname{span}\{W(e):e\in E_p^{\rm thick}\}
  =
  \operatorname{End}(\mathcal H_p^{\rm thick}).
\]
Finally,
\[
  \dim_\mathbb C\mathcal H_p^{\rm thick}
  =
  \prod_i |A_i|
  =
  \prod_i 3^{e_i\deg\pi_i}
  =
  3^{\deg p}
  =
  |R_p|.
\]
This is the sense in which thickenings add jet degrees of freedom at existing
closed points.

\subsection{Conclusion}

Nilpotent thickenings should manifest as local jet degrees of freedom.  They
do not create additional Zariski points.  The reduced theory records point
values over residue fields; the thickened theory quantizes the local Artin
rings and recovers the full coordinate-ring size.  The phase character detects
the socle, so the jet filtration is symplectically paired from the outside in:
order \(r\) position jets pair with order \(e-1-r\) momentum jets.
